\documentclass[11pt]{article}

% -------- Basics
\usepackage[margin=1in]{geometry}
\usepackage{amsmath,amssymb,mathtools}
\usepackage{hyperref}
\usepackage{enumitem}
\usepackage{graphicx}
\usepackage{booktabs}
\usepackage{siunitx}
\usepackage{microtype}
\usepackage{float} % for [H] figure placement
\usepackage[nameinlink,capitalise]{cleveref}

\setlist[itemize]{leftmargin=1.2em}
\setlist[enumerate]{leftmargin=1.2em}

% -------- Shortcuts
\newcommand{\E}{\mathbb{E}}
\newcommand{\Var}{\mathrm{Var}}
\newcommand{\Cov}{\mathrm{Cov}}
\newcommand{\softmax}{\operatorname{softmax}}
\newcommand{\logsumexp}{\operatorname{logsumexp}}
\newcommand{\R}{\mathbb{R}}

% -------- A tiny helper for dropping figures where you want them
% Usage: \inlinefig{width}{path}{caption}{label}
\newcommand{\inlinefig}[4]{%
  \begin{figure}[H]
    \centering
    \includegraphics[width=#1]{#2}
    \caption{#3}
    \label{#4}
  \end{figure}
}

\title{\textbf{KV merging}}
\author{Lerchen}
\date{\today}

\begin{document}
\maketitle

\section*{Motivation in one page}
We build a cache that respects two simple truths about attention:
\begin{enumerate}
  \item \textbf{Only relative key displacements matter.} For a fixed query $q$, adding a constant $c$ to \emph{all} keys shifts every logit by $q^\top c$, which cancels inside the softmax. So logits depend on $q^\top (k-\mu_k)$, not the absolute position of $k$. 
  \item \textbf{Queries are anisotropic.} Some directions have much larger second moment $u^\top\Sigma_q u$ than others
(eigenvalue spread can exceed $10\times$ in our Llama dumps).
\end{enumerate}
We recenter keys, whiten queries, and only merge when we can \emph{certify} an $L_2$ error bound under the empirical query distribution. 

\noindent In particular, for two keys \( k_1, k_2 \), define the discrepancy you actually care about, the squared difference in attention scores:
\[
D(k_1, k_2) := \mathbb{E}_q \big[ (q^\top (k_1 - k_2))^2 \big].
\]

\noindent However, note that 
\[D(k_1,k_2)=(k_1-k_2)^\top\,\Sigma_q\,(k_1-k_2)
=\big\|\Sigma_q^{1/2}(k_1-k_2)\big\|_2^2.\]

\noindent Hence, the correct representation of keys that clusters keys in a way that is connected with the \textit{difference} of the attention scores is exactly the whitened key vectors $\tilde k = \E[qq^\top]^{1/2} (k-\mu_k)$. We hence define formally:

\noindent \textbf{Definition.} Let $\Sigma_q=\E[qq^\top]$ (second moment, not the covariance) and let $\mu_k=\E[k]$ within a head. Define
\[
\tilde q \;=\; \Sigma_q^{-1/2}q,\qquad 
\tilde k \;=\; \Sigma_q^{1/2}(k-\mu_k).
\]
Then $\E[\tilde q \tilde q^\top]=I$ and
\[
\tilde q^\top \tilde k \;=\; q^\top (k-\mu_k).
\]
So attention logits can be written entirely in terms of the whitened query and the centered, reweighted key. 


\section*{Related work, why we differ}
KVMerger clusters in the raw space using cosine similarity and merges with a Gaussian weighting once a trigger fires.\footnote{KVMerger: \url{https://arxiv.org/abs/2407.08454}.} It assumes raw directions and absolute positions of keys are the right geometry, and it decides merges with heuristics, then checks quality after decoding. 
Q-Filters identifies a single query-driven direction and prunes by a projection score without materializing attention.\footnote{Q-Filters: \url{https://arxiv.org/abs/2503.02812}.} 

Our take:
\begin{itemize}
  \item Q-filters shows that the first right-vector of the SVD of $Q$ is a very dominant dimension, and that we should evict $k$'s that have smallest dot product with that direction. They justify this by saying that the \textit{keys have high magnitude in that direction}. This is a bad justification - one would want the \textit{variance} of the keys to be high in that dimension to be able to evict the odd $k$'s. We also generalize beyond only considering the right vectors of $Q$ as the dimensions - to find the dimensions that explain the \textit{most variance}, we are doing the best we can do - see next section.
  \item We choose to fit the quantity we actually care about. For a chunk of keys, we directly fit the merged key so that $\sum e^{q^\top k}$ is preserved in expectation, under the empirical $q$ distribution, and we accept a merge only if its certified error stays below $\epsilon$. 
  \item KVMerger uses cosine similarity between absolute positions of $k$ as a "similarity" metric many times. However, unlike our metric, this similarity metric does not respect shifts of $k$ and hence there is no reason for this to theoretically work - one could shift all of the $k$'s very very far in one direction, preserving all attention outputs yet making all $k$'s look very "similar", and this paradigm would elect to merge all of the $k$'s together into one $k$, a clear 
\end{itemize}

\subsection*{Why the spectrum of $S$ is the right guide.}

Let 
\[
S \;=\; \Cov(\tilde k)\;=\;\Sigma_q^{1/2}\,\Cov(k)\,\Sigma_q^{1/2}.
\]
Consider the centered logit $L(q,k)=\tilde q^\top \tilde k$. Its expected squared magnitude is
\[
\E_{q,k}\big[L(q,k)^2\big]
= \E_q\!\left[\tilde q^\top\,\E_k[\tilde k\tilde k^\top]\,\tilde q\right]
= \mathrm{Tr}\!\big(S\,\E_q[\tilde q \tilde q^\top]\big)
= \mathrm{Tr}(S).
\]
Projecting $\tilde k$ onto a rank-$r$ subspace with projector $P_r$ captures energy 
$\E\big[(\tilde q^\top P_r\tilde k)^2\big]=\mathrm{Tr}(P_r S)$. By Ky Fan, choosing the top-$r$ eigenvectors of $S$ maximizes this captured energy; it explains the most variance in attention scores. Empirically the spectrum drops fast: mean $\lambda_1=1.48\times 10^5$ falls to $3.96\times 10^4$ at rank $2$, so most useful logit energy sits in a tiny subspace once we make queries isotropic. We also hence motivate projecting all keys into a lower dimensional subspace consisting of the largest eigenvectors in $S$, via a projector $P_r$.

\subsection*{Learning the projector (optional but helps)}
We initialize $P_r$ with the top-$r$ eigenvectors of $S$ and then directly minimize the attention reconstruction loss
\[
\mathcal{L}_{\text{proj}} \;=\; \frac{1}{|B|}\!
\sum_{(Q,K,V,O)\in B}\big\|\mathrm{Attn}\!\big(\tilde Q,\,P_r\,\tilde K,\,V\big)-O\big\|_2^2,
\quad \tilde Q=\Sigma_q^{-1/2}Q,\qquad \tilde K=\Sigma_q^{1/2}(K-\mu_k)
\]
On cached DeepSeek batches, $\mathcal{L}_{\text{proj}}$ drops from $3.97\!\times\!10^{-3}$ at $r{=}1$ to $3.89\!\times\!10^{-4}$ at $r{=}64$, while a same-capacity per-head MLP trained on the identical loss stalls at $1.49\!\times\!10^{-2}$. The whitened linear subspace is both expressive and easy to optimize.

\section*{Least-squares merge certificate}
We merge only when we can certify low distortion under the empirical query distribution. For readability we omit the $\sqrt{d_k}$ terms.

\paragraph{Fitting the merged key (average-preserving).}
For a temporally contiguous chunk $\mathcal{C}$ and previously seen whitened queries $\{\tilde q_i\}_{i=1}^m$, define the average-of-exponentials target
\[
A_i \;=\; \frac{1}{|\mathcal{C}|}\sum_{k\in\mathcal{C}} e^{\tilde q_i^\top \tilde k}
\quad\Longrightarrow\quad
y_i \;=\; \log A_i \;=\; \log\!\sum_{k\in\mathcal{C}} e^{\tilde q_i^\top \tilde k} \;-\; \log|\mathcal{C}|.
\]
Stack the whitened queries into
\[
\mathbf{Q} \;=\;
\begin{bmatrix}
\tilde q_1^\top\\[-1pt]\vdots\\[-1pt]\tilde q_m^\top
\end{bmatrix}.
\]
We choose the merged whitened key by the (log-domain) least-squares fit
\[
\tilde k_\mathcal{C}
\;=\;\arg\min_{\tilde k}\big\|\mathbf{Q}\tilde k - y\big\|_2^2
\;=\;(\mathbf{Q}^\top\!\mathbf{Q})^{\dagger}\mathbf{Q}^\top y.
\]

\noindent \textit{Remark.} One can approximate the $(\mathbf{Q}^\top\!\mathbf{Q})^{\dagger}$ with $\frac 1m I$ in practice, as $\E [\tilde q \tilde q^t] = I$. 
\paragraph{Certifying distortion (original exp domain).}
We certify using the per-query, per-token squared error of the average-preserving approximation:
\[
\varepsilon_\mathcal{C} \;=\; \frac{1}{m\,|\mathcal{C}|}\sum_{i=1}^m 
\left(\sum_{k\in\mathcal{C}} e^{\tilde q_i^\top \tilde k}
\;-\; |\mathcal{C}|\, e^{\tilde q_i^\top \tilde k_\mathcal{C}}\right)^{\!2},
\]
and accept the merge only if $\varepsilon_\mathcal{C}\le\epsilon$.

\section*{Algorithm sketch}
\begin{enumerate}
  \item \textbf{Prefill.} Store singleton chunks $(\tilde k, v, 1)$.
  \item \textbf{Decode at step $t$:}
  \begin{enumerate}[label*=\arabic*.]
    \item Project the new key: $\tilde k_t \!=\! P_r\,\Sigma_q^{1/2}(k_t-\mu_k)$, append to the active chunk.
    \item Recompute $\tilde k_\mathcal{C}$ using the closed form above with
          $y_i=\log\!\sum_{k\in\mathcal{C}} e^{\tilde q_i^\top \tilde k}-\log|\mathcal C|$;
          set $v_\mathcal{C}$ to the mean of the values in the chunk.
    \item If $\varepsilon_\mathcal{C}>\epsilon$, seal the previous chunk and start a new one with $(\tilde k_t,v_t,1)$.
  \end{enumerate}
  \item \textbf{Lookup.} Each sealed chunk contributes a single term. We approximate
  \[
  \sum_{k\in\mathcal{C}} e^{q^\top k}\,V_k 
  \;\approx\;
  |\mathcal C|\,e^{\tilde q^\top \tilde k_\mathcal{C}}\cdot \bar v_\mathcal{C},
  \qquad
  \sum_{k\in\mathcal{C}} e^{q^\top k}
  \;\approx\;
  |\mathcal C|\,e^{\tilde q^\top \tilde k_\mathcal{C}},
  \]
  so keys and values share the same compression factor. Each chunk stores $(\tilde k_\mathcal{C},\,\bar v_\mathcal{C},\,|\mathcal{C}|)$.
\end{enumerate}

\section*{Snapshot of current evidence}
Using DeepSeek-R1 dumps:
\begin{itemize}
  \item In practice, it seems that actually using:

  \[
  \sum_{k\in\mathcal{C}} e^{q^\top k}\,V_k 
  \;\to\;
  \,e^{\tilde q^\top \tilde k_\mathcal{C}}\cdot \bar v_\mathcal{C},
  \qquad
  \sum_{k\in\mathcal{C}} e^{q^\top k}
  \;\to\;
  \,e^{\tilde q^\top \tilde k_\mathcal{C}},
  \]

  works better than scaling by these by $|\mathcal C|$ (i.e. when you do attention, you literally just have a cache of . This is surprising as this now does not resemble attention and does not attempt to get close to the numerator or denominator.
\end{itemize}

\subsection*{Appendix: AI use}
AI generated most of the code written in the Github repository, as well as most of the latex in the first edition of this document.

\noindent I had to fix many implementation bugs where the math got into the weeds, such as the least squares calculation and the low rank projection matmuls. I also had to replace much of the motivation and some of the math in this document with my own language; I hope that it reads more naturally and intuitively after this edit. Much "intuition" blabber in the first document was quite wrong.
\end{document}

